Este capítulo detalla el proceso de construcción del sistema \textbf{DOMU}, describiendo la estructura técnica del código, la configuración del entorno y las fases de implementación ejecutadas bajo la metodología híbrida definida. Se aborda la evolución del software desde su núcleo hasta los módulos de gestión comunitaria actualmente operativos.

\section{Estructura del Código y Arquitectura}
\label{sec:estructura_codigo}

El sistema se ha desarrollado siguiendo una arquitectura de cliente-servidor desacoplada, donde el Backend expone una API REST consumida por el Frontend. A continuación, se detalla la organización de los componentes en los repositorios de control de versiones.

\subsection{Backend (Java)}
El desarrollo del lado del servidor se ha estructurado en paquetes que separan claramente las responsabilidades, facilitando la mantenibilidad y escalabilidad del código. La estructura principal en el directorio \texttt{src/main/java/com/domu} es la siguiente:

\begin{itemize}
    \item \textbf{config/}: Clases de configuración global del sistema, incluyendo la inyección de dependencias (\texttt{AppConfig}, \texttt{DependencyInjectionModule}).
    \item \textbf{database/}: Capa de persistencia encargada de la comunicación directa con la base de datos MySQL (ej. \texttt{UserRepository}, \texttt{DataSourceFactory}).
    \item \textbf{domain/}: Entidades del dominio que representan los objetos de negocio como \texttt{User}, \texttt{Community}, \texttt{Ticket}, \texttt{Reservation} y \texttt{Vote}.
    \item \textbf{dto/}: Objetos de Transferencia de Datos (DTOs) utilizados para recibir y enviar información a través de la API sin exponer las entidades de la base de datos (ej. \texttt{LoginRequest}, \texttt{AuthResponse}, \texttt{VisitRequest}).
    \item \textbf{security/}: Módulos encargados de la seguridad, incluyendo la generación y validación de tokens JWT (\texttt{JwtProvider}) y el hashing de contraseñas (\texttt{BCryptPasswordHasher}).
    \item \textbf{service/}: Capa de lógica de negocio donde residen las reglas del sistema (ej. \texttt{UserService}).
    \item \textbf{web/}: Controladores y manejadores de rutas HTTP que exponen los endpoints de la API REST (\texttt{WebServer}, \texttt{UserMapper}).
\end{itemize}

\subsection{Frontend (React)}
La interfaz de usuario se ha construido utilizando React, organizando el código para promover la reutilización de componentes y la gestión eficiente del estado:

\begin{itemize}
    \item \textbf{components/}: Componentes reutilizables de la interfaz (botones, formularios, tarjetas de información).
    \item \textbf{pages/}: Vistas principales que componen la navegación del usuario (Login, Dashboard, Reservas, Votaciones).
    \item \textbf{services/}: Módulos encargados de realizar las peticiones HTTP al Backend (Axios/Fetch).
    \item \textbf{context/}: Manejo del estado global de la aplicación, como la sesión del usuario autenticado.
    \item \textbf{assets/}: Recursos estáticos como imágenes, logotipos e iconos.
\end{itemize}

\section{Fases de Implementación (Iteraciones)}
\label{sec:fases_implementacion}

Siguiendo la metodología híbrida adoptada, el desarrollo se dividió en fases iterativas, permitiendo desplegar funcionalidades incrementales. A continuación, se describen las iteraciones completadas y las que se encuentran en curso.

\subsection{Fase 1: Fundamentos y Seguridad}
El objetivo de esta primera iteración fue establecer los cimientos del sistema y asegurar el acceso controlado.
\begin{itemize}
    \item \textbf{Configuración del Entorno:} Levantar el servidor Backend y conectar con la base de datos MySQL.
    \item \textbf{Autenticación (JWT):} Se implementó un sistema de seguridad basado en \textit{JSON Web Tokens}. Al iniciar sesión, el usuario recibe un token con una duración de 1 hora, el cual se almacena en el \texttt{localStorage} del navegador para persistir la sesión durante su uso.
    \item \textbf{Gestión de Perfiles:} Funcionalidad para que los usuarios puedan editar su propia información de perfil.
\end{itemize}

\subsection{Fase 2: Gestión de Comunidades y Usuarios}
Se desarrollaron las herramientas administrativas necesarias para la operación del condominio.
\begin{itemize}
    \item \textbf{Registro de Comunidades:} Implementación de un flujo mediante correo electrónico para solicitar y validar el registro de nuevas comunidades en la plataforma.
    \item \textbf{Creación de Roles Operativos:} Desarrollo de interfaces para que el Administrador pueda crear y gestionar cuentas de personal, específicamente Conserjes y Personal de Aseo.
    \item \textbf{Vistas por Rol:} Segmentación de la interfaz para mostrar opciones específicas según si el usuario es Administrador, Conserje o Residente.
\end{itemize}

\subsection{Fase 3: Servicios al Residente y Operación}
Esta fase se centró en las funcionalidades de interacción diaria dentro de la comunidad.
\begin{itemize}
    \item \textbf{Reporte de Incidentes:} Módulo que permite a los residentes reportar problemas (ej. luces quemadas, ruidos molestos), generando un ticket visible para la administración.
    \item \textbf{Reservas de Áreas Comunes:} Sistema de calendario para verificar disponibilidad y reservar espacios como quinchos o salones.
    \item \textbf{Sistema de Votaciones:} Módulo digital que permite a la administración crear votaciones y a los residentes emitir su sufragio de manera remota.
\end{itemize}

\subsection{Fase 4: Control de Acceso y Finanzas (En Desarrollo)}
Actualmente, el proyecto se encuentra en esta etapa, enfocándose en los módulos más complejos de lógica de negocio.
\begin{itemize}
    \item \textbf{Gestión de Visitas (Formulario Manual):} Se ha implementado el formulario base que permite al conserje o al residente registrar los datos de una visita manualmente.
    \item \textbf{Generación de Gastos Comunes:} Desarrollo de la lógica para calcular el prorrateo de gastos mensuales por unidad.
    \item \textbf{Generación de Documentos:} Implementación de librerías para generar automáticamente los reportes de cobro en formato PDF.
\end{itemize}

\section{Detalles Técnicos de Implementación}

\subsection{Integración del Lector de Cédula (Diseño Técnico)}
Para el control de acceso en conserjería, se ha definido una estrategia de integración mediante hardware tipo "pistola escáner" de códigos QR/Barras. A diferencia de soluciones basadas en cámaras web, este dispositivo operará en modo de emulación de teclado (\textit{HID Mode}).

El flujo técnico diseñado (a implementar sobre el formulario actual) es el siguiente:
\begin{enumerate}
    \item El conserje enfoca el campo de entrada "RUT/Pasaporte" en la interfaz web.
    \item Al escanear el código QR de la Cédula de Identidad chilena, el dispositivo "escribe" automáticamente la cadena de texto contenida en el código.
    \item El Frontend procesará esta cadena, limpiando la URL o metadatos extraños para extraer únicamente el RUN.
    \item Se realizará una consulta automática a la base de datos (Endpoint \texttt{GET /api/visits/check/\{run\}}) para verificar si la persona tiene un acceso autorizado o histórico previo, autocompletando el formulario en caso positivo.
\end{enumerate}

\subsection{Manejo de Seguridad y Sesión}
La seguridad se gestiona mediante la clase \texttt{JwtProvider} en el Backend. El flujo de autenticación valida las credenciales contra la base de datos (con contraseñas encriptadas vía \texttt{BCrypt}) y emite un token firmado.

En el Frontend, se ha implementado un interceptor en las peticiones HTTP (servicio \textit{Axios/Fetch}) que inyecta automáticamente el token en la cabecera \texttt{Authorization: Bearer <token>}. Si el servidor responde con un error 401 (No autorizado) debido a la expiración del token (1 hora), el Frontend cierra la sesión automáticamente y redirige al Login por seguridad.

\section{Conclusión del Capítulo}
A la fecha, el sistema \textbf{DOMU} cuenta con un núcleo funcional sólido y los módulos de gestión comunitaria más críticos (votaciones, reservas, incidentes) totalmente operativos. El desarrollo actual se centra en finalizar el módulo financiero y la integración de hardware para optimizar el control de acceso, cumpliendo con el cronograma de la metodología híbrida establecida.